\documentclass[11pt]{article}
\usepackage[margin=0.8in]{geometry}
\usepackage{amsmath,microtype,hyperref}
\newcommand{\guideheading}[1]{\par\medskip\noindent{\large\bfseries #1}\par\smallskip}
\begin{document}
\begin{center}
{\Large Guide: A Variable-Length Route Past Some Table Gaps}\par\medskip
Collatz proof project\quad---\quad September 5, 2026
\end{center}
\guideheading{Why this matters}
The finite table can fail throughout an arbitrarily long initial odd run.
That does not mean every merging argument fails there. A classical
Garner rule has a length determined by the run itself. We have now
formalized its arithmetic applicability and the exact case it leaves open.
These identities are known mathematics, also treated explicitly by
\href{https://arxiv.org/abs/1709.02979}{LaDue (2017)}.

\guideheading{The two exits}
Write an odd seed as $n=2^am-1$, where $m$ is odd. The number $a$ is
the length of the initial odd run. After that run the orbit takes an
even step. What happens next determines the early comparison with $n-1$.

If the next step is also even, the two seeds reach the same value after
$a+2$ steps. This works for every run length. If the next step is odd,
their values at time $a+2$ are instead
\[
9z+2\quad\text{and}\quad z,
\]
where $z$ is the value reached from $n-1$. Lean checks both branches
and the precise arithmetic condition selecting them.

\guideheading{Which minimality argument applies}
Suppose $n>1$ were the least positive seed with an unbounded orbit.
The smaller seed $n-1$ would have a bounded orbit. If the two seeds
merge, their shared tail makes $n$ bounded too, a contradiction.

This argument needs ordinary merging. It does not require the smaller
seed to have noncontracting coefficients. In fact $n-1$ is even, so its
first coefficient is below one. Requiring every replacement to preserve
full noncontraction would miss this useful classical rule.

\guideheading{Connection to the previous obstruction}
We can construct arbitrarily large members of the table-obstruction
family that take the double-even exit. The table only retraces their
early odd steps, but the variable-length Garner rule merges them with
their predecessor later. We can also construct members taking the
single-even exit. The obstruction therefore contains both cases;
the new application does not dispose of the whole family.

\guideheading{The remaining gap}
In the single-even case, $z$ lies on the bounded orbit of $n-1$,
while $9z+2$ lies on the supposedly unbounded orbit of $n$.
We have not proved a rule that transfers boundedness from $z$ to $9z+2$.
The exact relationship identifies the unresolved problem; it does not
solve it. Collatz remains unproved, including the nontrivial-cycle issue.

\guideheading{How to verify the work}
\begingroup\raggedright
The formal source is \texttt{lean/Collatz/OddRunMerges.lean}.
From \texttt{lean/}, run \texttt{lake build Collatz.OddRunMerges}.
From the root, run \texttt{python3 analysis/odd\_run\_merges.py} for
constructed examples or \texttt{python3 -m unittest discover -s analysis
-p test\_odd\_run\_merges.py} for independent orbit checks.
The generic endpoint identities and least-unbounded exclusion are
kernel checked. The CRT application to the obstruction is a written
deduction, supported by exact instance checks.
\par\endgroup
\end{document}
